% GENERATED FILE. Edit canonical lessons, then run npm run generate:books.
% canonical-source-hash: fnv1a64:c671d9f8bb51ff39
% canonical-lessons: ZH-C274-huabing, ZH-C274-shouji, ZH-C274-remen, ZH-C274-shougong, ZH-C274-qiumi

\chapter{Free Time}
\label{ch:zh-free-time}
\hlchaptermodality{274}

\begin{chapteropening}
\textbf{By the end of this chapter:} \emph{I can talk about skating, collecting things, what is all the rage, handicrafts and sports fans.}

Complete the last lesson of chapter 274: I can talk about skating, collecting things, what is all the rage, handicrafts and sports fans.
\end{chapteropening}

\section[huábīng]{huábīng — to ice-skate}

\label{lesson:ZH-C274-huabing}



\begin{quote}
Before the new one: say the Mandarin for careless, then the Mandarin for careful with details.
\end{quote}

\subsection*{You'll want to know: huábīng}
\emph{huábīng} — \textquotedblleft{}to ice-skate\textquotedblright{}.

\begin{quote}
\emph{dōngtiān wǒmen qù huábīng} — In winter we go skating.
\end{quote}

Its characters come later; for now it is a word for the ear.

\begin{cousinweb}
\emph{huábīng} is built from \emph{huá} 'slippery' and \emph{bīng} 'ice', both inherited from Middle Chinese, \emph{hweat} and \emph{ping} (Guangyun, in Baxter's transcription), which glosses the first 'smooth' and quotes an older dictionary for the second, 'frozen water'.
\end{cousinweb}

\begin{grammarlens}[title={What you've built}]
One more word for what I like.
\end{grammarlens}

\subsection*{Your turn}
\begin{itemize}
\raggedright
  \item \emph{Say it:} \emph{huábīng}
  \item \emph{Say it:} \emph{huábīng}, once more
  \item \emph{Say it:} \emph{dōngtiān wǒmen qù huábīng} — In winter we go skating.
  \item \emph{Recall:} say the Mandarin for careless, then the Mandarin for careful with details, then say \emph{huábīng} again
\end{itemize}

\begin{tcolorbox}[breakable,colback=teal!4,colframe=teal!35!black,title={Before you move on}]
What does huábīng mean? (\textquotedblleft{}To ice-skate\textquotedblright{}.) Say it once more.
\end{tcolorbox}
\section[shōují]{shōují — to collect things}

\label{lesson:ZH-C274-shouji}



\begin{quote}
Before the new one: say the Mandarin for careful with details, then the Mandarin for to ice-skate.
\end{quote}

\subsection*{You'll want to know: shōují}
\emph{shōují} — \textquotedblleft{}to collect things\textquotedblright{}.

\begin{quote}
\emph{wǒ xǐhuan shōují yóupiào} — I like collecting stamps.
\end{quote}

Its characters come later; for now it is a word for the ear.

\begin{cousinweb}
\emph{shōují} is built from \emph{shōu} 'to gather in' and \emph{jí} 'to assemble', both inherited from Middle Chinese, \emph{syuw} and \emph{dzip} (Guangyun, in Baxter's transcription), which glosses the first 'to gather in' and the second 'to gather'.
\end{cousinweb}

\begin{grammarlens}[title={What you've built}]
One more word for what I like.
\end{grammarlens}

\subsection*{Your turn}
\begin{itemize}
\raggedright
  \item \emph{Say it:} \emph{shōují}
  \item \emph{Say it:} \emph{shōují}, once more
  \item \emph{Say it:} \emph{wǒ xǐhuan shōují yóupiào} — I like collecting stamps.
  \item \emph{Recall:} say the Mandarin for careful with details, then the Mandarin for to ice-skate, then say \emph{shōují} again
\end{itemize}

\begin{tcolorbox}[breakable,colback=teal!4,colframe=teal!35!black,title={Before you move on}]
What does shōují mean? (\textquotedblleft{}To collect things\textquotedblright{}.) Say it once more.
\end{tcolorbox}
\section[\texorpdfstring{\zh{热门} (rèmén)}{热门 (rèmén)}]{\zh{热门} (rèmén) — popular; all the rage}

\label{lesson:ZH-C274-remen}



\begin{quote}
Before the new one: say the Mandarin for to ice-skate, then the Mandarin for to collect things.
\end{quote}

\subsection*{You'll want to know: \zh{热门}}
\textbf{\zh{热门}} — \emph{rèmén} — \textquotedblleft{}popular; all the rage\textquotedblright{}.

\begin{quote}
\textbf{\zh{这本小说很热门}} \emph{zhè běn xiǎoshuō hěn rèmén} — This novel is all the rage.
\end{quote}

\begin{cousinweb}
\textbf{\zh{热门}} (\emph{rèmén}) is built from \zh{热} 'hot' and \zh{门} 'a door', both inherited from Middle Chinese, \emph{nyet} and \emph{mwon} (Guangyun, in Baxter's transcription). How a hot door came to mean something popular is not settled.
\end{cousinweb}

\begin{grammarlens}[title={What you've built}]
One more word for what I like.
\end{grammarlens}

\subsection*{Your turn}
\begin{itemize}
\raggedright
  \item \emph{Say it:} \emph{rèmén}
  \item \emph{Say it:} \emph{rèmén}, once more
  \item \emph{Say it:} \emph{zhè běn xiǎoshuō hěn rèmén} — This novel is all the rage.
  \item \emph{Recall:} say the Mandarin for to ice-skate, then the Mandarin for to collect things, then say \emph{rèmén} again
\end{itemize}

\begin{tcolorbox}[breakable,colback=teal!4,colframe=teal!35!black,title={Before you move on}]
What does \zh{热门} mean? (\textquotedblleft{}Popular; all the rage\textquotedblright{}.) Say it once more.
\end{tcolorbox}
\section[\texorpdfstring{\zh{手工} (shǒugōng)}{手工 (shǒugōng)}]{\zh{手工} (shǒugōng) — handicraft; things made by hand}

\label{lesson:ZH-C274-shougong}



\begin{quote}
Before the new one: say the Mandarin for to collect things, then the Mandarin for all the rage.
\end{quote}

\subsection*{You'll want to know: \zh{手工}}
\textbf{\zh{手工}} — \emph{shǒugōng} — \textquotedblleft{}handicraft; things made by hand\textquotedblright{}.

\begin{quote}
\textbf{\zh{我们先做手工}} \emph{wǒmen xiān zuò shǒugōng} — Let's do the crafts first.
\end{quote}

\begin{cousinweb}
\textbf{\zh{手工}} (\emph{shǒugōng}) is built from \zh{手} 'the hand' and \zh{工} 'work', both inherited from Middle Chinese, \emph{syuwX} and \emph{kuwng} (Guangyun, in Baxter's transcription).
\end{cousinweb}

\begin{grammarlens}[title={What you've built}]
One more word for what I like.
\end{grammarlens}

\subsection*{Your turn}
\begin{itemize}
\raggedright
  \item \emph{Say it:} \emph{shǒugōng}
  \item \emph{Say it:} \emph{shǒugōng}, once more
  \item \emph{Say it:} \emph{wǒmen xiān zuò shǒugōng} — Let's do the crafts first.
  \item \emph{Recall:} say the Mandarin for to collect things, then the Mandarin for all the rage, then say \emph{shǒugōng} again
\end{itemize}

\begin{tcolorbox}[breakable,colback=teal!4,colframe=teal!35!black,title={Before you move on}]
What does \zh{手工} mean? (\textquotedblleft{}Handicraft; things made by hand\textquotedblright{}.) Say it once more.
\end{tcolorbox}
\section[qiúmí]{qiúmí — sports fan}

\label{lesson:ZH-C274-qiumi}



\begin{quote}
Before the new one: say the Mandarin for all the rage, then the Mandarin for handicrafts.
\end{quote}

\subsection*{You'll want to know: qiúmí}
\emph{qiúmí} — \textquotedblleft{}sports fan\textquotedblright{}.

\begin{quote}
\emph{wǒ bàba shì yí ge qiúmí} — My dad is a big sports fan.
\end{quote}

Its characters come later; for now it is a word for the ear.

\begin{cousinweb}
\emph{qiúmí} is built from \emph{qiú} 'a ball' and \emph{mí}, both inherited from Middle Chinese, \emph{gjuw} and \emph{mej} (Guangyun, in Baxter's transcription), which glosses the second 'to be led astray'.
\end{cousinweb}

\begin{grammarlens}[title={What you've built}]
One more word for what I like.
\end{grammarlens}

\subsection*{Your turn}
\begin{itemize}
\raggedright
  \item \emph{Say it:} \emph{qiúmí}
  \item \emph{Say it:} \emph{qiúmí}, once more
  \item \emph{Say it:} \emph{wǒ bàba shì yí ge qiúmí} — My dad is a big sports fan.
  \item \emph{Recall:} say the Mandarin for all the rage, then the Mandarin for handicrafts, then say \emph{qiúmí} again
\end{itemize}

\begin{tcolorbox}[breakable,colback=teal!4,colframe=teal!35!black,title={Before you move on}]
What does qiúmí mean? (\textquotedblleft{}Sports fan\textquotedblright{}.) Say it once more.
\end{tcolorbox}
